\section{Discussion, Boundaries, and Related Work}
\label{sec:discussion}

\paragraph{Comparison.}
Against model-level watermarks (KGW and descendants; SynthID-Text), AAWM
trades sampling-level control for black-box applicability and per-user
attribution; against PostMark-class post-hoc watermarks, AAWM provides
multi-user uid decoding (16 bits with registry matching and abstention)
where presence-ratio schemes carry essentially one bit of existence;
against cryptographic agent identity, AAWM marks the text itself rather
than the agent's credentials. The honest price of uid capacity is
exposure to the pitfalls we characterize in the companion paper:
correlated-vote inflation (mitigated by per-word caps), domain-specific
null models (mitigated by in-domain calibration), and template-class
adversarial texts (mitigated by repetition routing)---all three mitigations
are shipped in v0.15.0.

\paragraph{Boundaries.}
(1)~Semantic LLM rewriting destroys the mark (Table~\ref{tab:ea2}).
(2)~The English zero-cost dictionary is sparse on non-AI text; production
English deployments should use the full-lexicon codec.
(3)~Legal-grade attribution is out of scope: the mark is statistical, not
cryptographic evidence.
(4)~A symmetric master key means verifier compromise is total; key
management is an operational precondition.

\paragraph{Artifact.}
AAWM is open source (MIT core) at v0.15.0: Python SDK, CLI, HTTP service,
five framework adapters, fail-open middleware, key rotation, audit
logging, and the full experiment suite of this paper with fixed seeds.
